\documentclass[11pt]{article}

\usepackage[margin=1in]{geometry}
\usepackage[hidelinks]{hyperref}

\title{Reliability modeling for transportable perturbation-response predictions in public single-cell functional-genomics screens}
\author{Yiheng Yang\\
School of Future Technology, Dalian University of Technology\\
School of Bioengineering, Dalian University of Technology\\
\texttt{yangyiheng@mail.dlut.edu.cn}}
\date{Prepared for single-anonymized submission to Computational Biology and Chemistry}

\begin{document}
\maketitle

\section*{Corresponding Author}
Yiheng Yang\\
School of Future Technology, Dalian University of Technology\\
School of Bioengineering, Dalian University of Technology\\
Email: \texttt{yangyiheng@mail.dlut.edu.cn}

\section*{Funding}
This research did not receive any specific grant from funding agencies in the public, commercial, or not-for-profit sectors.

\section*{Declaration of Competing Interest}
The author declares no competing interests.

\section*{Acknowledgements}
The author thanks the maintainers and contributors of the public single-cell perturbation data resources used in this study.

\section*{Data and Code Availability}
The minimal code release for peer review is available as a tagged source snapshot at \url{https://github.com/Neabigmo/PTL/tree/v1.0-cbac-submission}. The current evidence surface contains nine processed public screens and 525 baseline runs. Processed evidence tables, split manifests, PTL outputs, case-example tables, and figure source data are included with the submission package as a repository-deposition-ready archive. Raw public datasets are not redistributed and remain available from their original repositories.

\section*{CRediT Author Statement}
Yiheng Yang conceived the study, designed the benchmark, implemented the analysis workflow, interpreted the results, and prepared the manuscript.

\end{document}
